% lit-ext-proofs.tex -- cluster "lit-ext" (prior work for conditional recourse, constraints,
% optimal sets and structural limits).  Fragment without preamble.  Citation keys refer to
% lit-ext.bib.  Labels of the form sec:..., lem:round, prop:filter, lem:contraction,
% lem:localize, lem:states, thm:approx, thm:exact, eq:growth, eq:semiconcavity refer to the
% corresponding statements of the main paper (main.tex of the report).
%
% Contents
%   A. Classical facts about conditional values (with attribution).
%   B. Exact grid oracles for balanced quadratics: a width-free instance of the certified
%      grid algorithm (new corollary obtained in this audit).
%   C. Optimal sets: the face criterion and the diagonal Lagrangian certificate.
%   D. Mean-preserving rounding in totally unimodular fibres.
%   E. Draft of the related-work section.

\section*{A. Conditional values: concavity and curvature}

Throughout, a function $\phi$ on a box has \emph{upper coordinate curvature} $\ell_j$ in
coordinate $j$ if $t\mapsto\phi(v+te_j)-\ell_jt^2/2$ is concave on every coordinate segment
contained in the box. This is coordinatewise semiconcavity with a linear modulus in the sense
of \cite{CannarsaSinestrari2004}.

\begin{lemma}[Infima preserve concavity and coordinate curvature]\label{lit:lem:inf}
Let $V\subseteq\R^k$ be convex, let $\mathcal Y$ be a nonempty index set, and let
$\{f_y\}_{y\in\mathcal Y}$ be functions on $V$ with $\phi(v)=\inf_{y}f_y(v)>-\infty$ for all
$v\in V$.
\begin{enumerate}
\item If every $f_y$ has upper coordinate curvature $\ell_j$ in coordinate $j$, so does $\phi$.
\item If $f_y(v)=a(y)+w(y)^{\mathsf T}v+\tfrac12v^{\mathsf T}Av+d^{\mathsf T}v+e$ with $A$
symmetric and $A,d,e$ independent of $y$, then $\phi(v)-\tfrac12v^{\mathsf T}Av-d^{\mathsf T}v$
is concave on $V$. Consequently $\phi$ has upper coordinate curvature $\ell_j$ for every
$\ell_j\ge A_{jj}$.
\end{enumerate}
\end{lemma}

\begin{proof}
For (1), fix a coordinate segment and write $g_y(t)=f_y(v+te_j)-\ell_jt^2/2$, which is concave.
For $t_1,t_2$ in the segment and $\lambda\in[0,1]$,
\[
\inf_y g_y(\lambda t_1+(1-\lambda)t_2)\ge\inf_y\bigl[\lambda g_y(t_1)+(1-\lambda)g_y(t_2)\bigr]
\ge\lambda\inf_y g_y(t_1)+(1-\lambda)\inf_y g_y(t_2),
\]
and $\inf_yg_y(t)=\phi(v+te_j)-\ell_jt^2/2$. For (2), $f_y(v)-\tfrac12v^{\mathsf T}Av-d^{\mathsf T}v$
is affine in $v$; a finite pointwise infimum of affine functions is concave by the same
inequality. On a coordinate line the subtracted quadratic is $A_{jj}t^2/2$ plus an affine
function, so $\phi(v+te_j)-A_{jj}t^2/2$ is concave; subtracting $(\ell_j-A_{jj})t^2/2\ge0$ preserves
concavity.
\end{proof}

\begin{remark}[Use and attribution]
Part (2) gives, for a private convex block with a parameter-independent feasible set,
$\psi(v)=\min_{y\in Y}\{\tfrac12y^{\mathsf T}Cy+y^{\mathsf T}Dv+c^{\mathsf T}y\}$, that $\psi$ is
concave, so a value factor adds no positive coordinate curvature. Part (1) gives the conditional
value $W_B(v)=\min_{x_{-B}}F(v,x_{-B})$ the same curvature bound as $F$ when the outside domain
does not depend on $v$. Both are classical: concavity of optimal values in parameters that enter
the objective concavely over a fixed feasible set is recorded in \cite{MangasarianRosen1964} and
systematically in \cite{FiaccoKyparisis1986}; stability of semiconcavity under infima is standard
\cite{CannarsaSinestrari2004}. Neither the convexity of the inner problem nor uniqueness of the
inner minimizer is used.
\end{remark}

\begin{lemma}[Concave kinks]\label{lit:lem:kinks}
Let $g:[\alpha,\omega]\to\R$ be continuous and let $\alpha=\tau_0<\dots<\tau_N=\omega$. Suppose
that on each $[\tau_{r-1},\tau_r]$ the function $g$ agrees with a quadratic $q_r$ with
$q_r''\le\ell$, and that $g-\psi$ is concave for some continuously differentiable $\psi$. Then
$t\mapsto g(t)-\ell t^2/2$ is concave on $[\alpha,\omega]$.
\end{lemma}

\begin{proof}
Put $h=g-\ell t^2/2$. On each piece $h$ is a quadratic with $h''\le0$. A concave function has
left derivative at least its right derivative at every interior point; adding the
$C^1$ function $\psi-\ell t^2/2$ does not change the difference of one-sided derivatives. Hence
$h'_-(\tau_r)\ge h'_+(\tau_r)$ for $0<r<N$. The right derivative $h'_+$ therefore exists on
$[\alpha,\omega)$ and is nonincreasing: within a piece it is the derivative of a concave
quadratic, and across $\tau_r$ it satisfies $h'_+(t)\ge h'_-(\tau_r)\ge h'_+(\tau_r)$ for
$t<\tau_r$ in the previous piece. Since $h$ is continuous and piecewise $C^1$,
$h(t)=h(\alpha)+\int_\alpha^th'_+$, and a function with a nonincreasing derivative of this kind is
concave.
\end{proof}

\begin{corollary}[Curvature across a certified partition]\label{lit:cor:partition}
Let $f(y,v)=\tfrac12y^{\mathsf T}Cy+y^{\mathsf T}Dv+c^{\mathsf T}y+\tfrac12v^{\mathsf T}Av+d^{\mathsf T}v+e$
with $C\succeq0$, $y\in Y=[\underline y,\overline y]$ and $v$ in a box $V$. Suppose $V$ is covered by finitely many
closed convex polyhedra $P_r$ and on each $P_r$ the conditional value
$\phi(v)=\min_{y\in Y}f(y,v)$ equals a quadratic $q_r$ with Hessian $H_r$. Then $\phi$ has
upper coordinate curvature $\max_r(H_r)_{jj}$ in coordinate $j$.
\end{corollary}

\begin{proof}
The function $\phi$ is continuous because $f$ is continuous and $Y$ is compact. A
coordinate segment meets each $P_r$ in a closed interval. Discarding degenerate intervals and
refining overlaps (on an overlap of positive length two quadratics that both equal $\phi$ coincide)
yields consecutive intervals on which $\phi$ is a quadratic with second derivative
$(H_r)_{jj}$ in $t$. By Lemma~\ref{lit:lem:inf}(2) with $\mathcal Y=Y$, $\phi$ minus a
smooth quadratic is concave. Lemma~\ref{lit:lem:kinks}, applied with curvature bound $\max_r(H_r)_{jj}$, concludes.
\end{proof}

\begin{remark}
When each $P_r$ carries an affine primal response and affine nonnegative bound multipliers
satisfying the KKT conditions on $P_r$, the identity
$f(y,v)-f(\hat y_r(v),v)=\tfrac12(y-\hat y_r)^{\mathsf T}C(y-\hat y_r)+\lambda_r^{\mathsf T}(y-\underline y)
+\mu_r^{\mathsf T}(\overline y-y)\ge0$ certifies $\phi=q_r$ on $P_r$. For $C\succ0$ and linearly independent
active constraints such regions are the critical regions of multiparametric quadratic
programming \cite[Theorems~2 and~4]{BemporadEtAl2002}; for merely positive semidefinite $C$, where
optimizers can be nonunique, see \cite{SpjotvoldTondelJohansen2005,SpjotvoldTondelJohansen2007,
PatrinosSarimveis2011}. Corollary~\ref{lit:cor:partition} only needs the value function on each
region, so no global overlay of the regions of different blocks is required. The concave-kink
structure of univariate value functions is also the key structural fact of the forest algorithm
of \cite{DelPiaKhajavirad2026}.
\end{remark}

\begin{lemma}[Constant gradient on the optimal set of a convex quadratic]\label{lit:lem:mangasarian}
Let $K\subseteq\R^m$ be convex and $f(y)=\tfrac12y^{\mathsf T}Cy+c^{\mathsf T}y$ with $C\succeq0$.
If $S=\arg\min_Kf\neq\emptyset$ and $\bar y\in S$, then
\[
S=\{y\in K:\ Cy=C\bar y,\ c^{\mathsf T}y=c^{\mathsf T}\bar y\}.
\]
In particular all minimizers have the same gradient $C\bar y+c$.
\end{lemma}

\begin{proof}
For $y_1,y_2\in S$, $(y_1+y_2)/2\in K$ and the exact identity
$f(\tfrac{y_1+y_2}2)=\tfrac12f(y_1)+\tfrac12f(y_2)-\tfrac18(y_1-y_2)^{\mathsf T}C(y_1-y_2)$ gives
$(y_1-y_2)^{\mathsf T}C(y_1-y_2)\le0$, hence $C(y_1-y_2)=0$ because $C\succeq0$. Then
$y_1^{\mathsf T}Cy_1=y_1^{\mathsf T}Cy_2=y_2^{\mathsf T}Cy_2$, and $f(y_1)=f(y_2)$ forces
$c^{\mathsf T}y_1=c^{\mathsf T}y_2$. Conversely, if $y\in K$ satisfies the two displayed equations,
the same computation gives $f(y)=f(\bar y)$.
\end{proof}

\begin{remark}
This is \cite[Lemma~1 and Corollary~1]{Mangasarian1988}. It is the fact behind the observation
that all central minimizers of a private convex block share one gradient, which fixes the
sign pattern of every globally affine optimal selector.
\end{remark}

\section*{B. Exact grid oracles for balanced quadratics}

The certified grid algorithm uses its tree decomposition only to compute, on the current product
grid $\prod_iG_i$, the corrected minimum $\beta=\min Q$, one minimizer, and the coordinate
min-marginals $m_i(v)$ of $Q=F-D$. Any exact oracle for these quantities can replace the
messages. For one sign structure such an oracle exists without any width bound.

\begin{definition}[Balanced quadratic]
A quadratic $F(x)=\tfrac12x^{\mathsf T}Hx+b^{\mathsf T}x+e$ is \emph{balanced} if there are signs
$\sigma\in\{-1,1\}^n$ with $\sigma_i\sigma_jH_{ij}\le0$ for all $i\neq j$.
\end{definition}

Equivalently, after the substitution $x_i\mapsto\sigma_ix_i$ all off-diagonal Hessian entries
are nonpositive, that is, $F$ becomes submodular on $\R^n$. No sign condition is imposed on the
diagonal or the linear terms.

\begin{lemma}[Recognizing balance]\label{lit:lem:balance}
Balance can be decided, and signs $\sigma$ found, in time linear in $n$ plus the number of
nonzero off-diagonal entries.
\end{lemma}

\begin{proof}
The conditions are $\sigma_i=\sigma_j$ when $H_{ij}<0$ and $\sigma_i=-\sigma_j$ when $H_{ij}>0$.
Breadth-first search assigns signs on each connected component of the graph of nonzero entries
and reports failure exactly when an edge contradicts the assignment, that is, when some cycle has
an odd number of positive entries \cite{Harary1953}.
\end{proof}

\begin{lemma}[Grid minimization by one minimum cut]\label{lit:lem:chaincut}
Let $F$ be balanced with signs $\sigma$, let $G_i\subset\R$ be finite and nonempty with
$|G_i|=m_i+1$, and let $\eta_i:G_i\to\Q$ be arbitrary. Put
$\Phi(y)=F(y)+\sum_i\eta_i(y_i)$ for $y\in\prod_iG_i$. Then $\min\Phi$ and a minimizer can be
computed from one minimum $s$--$t$ cut in a directed graph with $\sum_im_i+2$ nodes and at most
$3\sum_im_i+2\sum_{i<j:\,H_{ij}\neq0}m_im_j$ arcs whose capacities are computed from the
coefficients, the grid points and the values $\eta_i$ by polynomially many rational operations.
A maximum flow of equal value certifies the minimum. The same
holds for each min-marginal $\min\{\Phi(y):y_i=v\}$, by replacing $G_i$ with $\{v\}$.
\end{lemma}

\begin{proof}
List $G_i$ as $g_{i,0},\dots,g_{i,m_i}$ in increasing order if $\sigma_i=1$ and decreasing order
if $\sigma_i=-1$, so that $\delta_{i,k}=g_{i,k}-g_{i,k-1}$ has the sign of $\sigma_i$. Encode
$y_i=g_{i,k}$ by $z_i\in\{0,1\}^{m_i}$ with $z_{i,l}=1$ exactly for $l\le k$. This is a bijection
between $G_i$ and the monotone vectors $z_{i,1}\ge\dots\ge z_{i,m_i}$, and on monotone vectors
\[
y_i=g_{i,0}+\sum_l\delta_{i,l}z_{i,l},\qquad
\theta(y_i)=\theta(g_{i,0})+\sum_l\bigl[\theta(g_{i,l})-\theta(g_{i,l-1})\bigr]z_{i,l}
\]
for every function $\theta$ on $G_i$. Apply the second formula to
$\theta_i(t)=\tfrac12H_{ii}t^2+b_it+\eta_i(t)$ and substitute the first into each cross term
$H_{ij}y_iy_j$, $i<j$. On monotone vectors this gives
\[
\Phi=E(z)=E_0+\sum_ah_az_a+\sum_{a<b}w_{ab}z_az_b,\qquad
w_{(i,k),(j,l)}=H_{ij}\delta_{i,k}\delta_{j,l}\le0,
\]
for suitable rational $E_0,h_a$; all other pair coefficients are zero. For binary
variables $z_az_b=\tfrac12(z_a+z_b)-\tfrac12|z_a-z_b|$, so with $r_{ab}=-w_{ab}/2\ge0$ and
$p_a=h_a+\tfrac12\sum_{b\neq a}w_{ab}$,
\[
E(z)=E_0+\sum_ap_az_a+\sum_{a<b}r_{ab}|z_a-z_b|.
\]
Build a graph with a node for each $z_a$, an arc $a\to t$ of capacity $\max\{p_a,0\}$, an arc
$s\to a$ of capacity $\max\{-p_a,0\}$, two opposite arcs of capacity $r_{ab}$ between $a$ and
$b$, and an arc $(i,l+1)\to(i,l)$ of infinite capacity for each $i$ and $1\le l<m_i$. Identify
$z$ with the cut whose source side is $\{s\}\cup\{a:z_a=1\}$. A finite cut has no
infinite arc from the source side to the sink side, which says exactly that $z$ is monotone; and
for monotone $z$ the cut capacity is
$\sum_a[\max\{p_a,0\}z_a+\max\{-p_a,0\}(1-z_a)]+\sum_{a<b}r_{ab}|z_a-z_b|$, so that
\[
E(z)=E_0+\sum_a\min\{p_a,0\}+\operatorname{cut}(z).
\]
The all-zero vector is monotone, so the minimum cut is finite and decodes to a minimizer of
$\Phi$. By max-flow min-cut duality, a feasible flow whose value equals the capacity of the
returned cut certifies its optimality. Fixing $y_i=v$ is the same construction with $G_i=\{v\}$,
and the restricted function is still balanced.
\end{proof}

\begin{remark}
Encoding finite ordered labels by threshold variables, and solving pairwise problems that are
submodular with respect to the label orders by one minimum cut, is classical; see
\cite[Section~4]{SchlesingerFlach2006} and, for convex functions of label differences,
\cite{Ishikawa2003}. For binary variables the construction is that of
\cite{Ivanescu1965,PicardRatliff1975,KolmogorovZabih2004}. The point of
Lemma~\ref{lit:lem:chaincut} here is only that the arbitrary unary terms $\eta_i$ may be the
negative interpolation corrections $-d_i$ on arbitrary nonuniform grids.
\end{remark}

\begin{theorem}[Balanced box quadratics without a width parameter]\label{lit:thm:balanced}
Let $F(x)=\tfrac12x^{\mathsf T}Hx+b^{\mathsf T}x+e$ be an explicit rational quadratic on a
mixed box $X$ with binary input length $I$, and suppose $F$ is balanced. Let
$L=\max_iH_{ii}$.
\begin{enumerate}
\item If $L\le0$, an exact minimizer and the optimal value are obtained from one minimum cut on
the endpoint labels, in time polynomial in $I$.
\item If $L>0$ and $F$ has a unique minimizer satisfying the growth condition~\eqref{eq:growth}
with constant $g>0$, put $\kappa=\max\{1,L/g\}$. The capped algorithm of
Section~\ref{sec:rate}, with $\beta$, a minimizer and all min-marginals computed by
Lemma~\ref{lit:lem:chaincut} instead of tree messages, returns a feasible rational point and an
independently checkable certificate of gap at most $2^{-q}$ in
$(n\kappa(I+q+1))^{O(1)}$ bit operations. The exact minimizer and optimal value are returned in
$(n\kappa(I+1))^{O(1)}$ bit operations.
\end{enumerate}
No decomposition is used, $g$ and $\kappa$ are not inputs, and the validity of every certificate
uses neither growth nor uniqueness.
\end{theorem}

\begin{proof}
(1) If every $H_{ii}\le0$, then $F$ is concave along every coordinate line of the continuous
hull, so successive replacement of coordinates by endpoints of their (inward-rounded) intervals
does not increase $F$, and the minimum over $X$ equals the minimum over the endpoint labels.
Lemma~\ref{lit:lem:chaincut} with $|G_i|\le2$ and $\eta_i=0$ computes it.

(2) The correctness statements used by the capped algorithm, namely Lemma~\ref{lem:round},
Proposition~\ref{prop:filter}, Lemma~\ref{lem:contraction}, Lemma~\ref{lem:localize} and
Lemma~\ref{lem:states}, depend on $F$ only through the coordinate curvature
condition~\eqref{eq:semiconcavity}, which holds with this $L$ because the restriction of $F$ to a
coordinate line has second derivative $H_{ii}$, and through the exact values $\beta$, a
minimizer of $Q$, and $m_i(v)$ on the current grids. The tree decomposition enters those proofs
only through the computation of these values by \eqref{eq:messages}--\eqref{eq:calibrated}.
Since $Q=F-D$ and $D$ is a sum of unary terms, Lemma~\ref{lit:lem:chaincut} with
$\eta_i=-d_i$ computes $\beta$ and a minimizer by one cut, and each $m_i(v)$ by one further cut.
Hence all correctness and state bounds hold verbatim, including the abort rule of the capped
trials, and the certificate consists of the same filtering record with each tree computation
replaced by a cut and an equal-value flow.

For the work, the trial with $\theta=2^{-\mu}$ has at most
$K_\mu=100\cdot2^\mu\lceil\log_2(n+2)\rceil$ nodes per coordinate (Lemma~\ref{lem:states}),
and the first successful trial has $2^\mu\le6\sqrt\kappa$. A stage therefore performs at most
$1+nK_\mu$ cut computations on graphs with at most $nK_\mu+2$ nodes and
$3nK_\mu+n^2K_\mu^2$ arcs. As in the proof of Theorem~\ref{thm:approx}, grid coordinates at
stage $j$ have denominators dividing $A_j=D2^{j+\mu K_\mu}$ and all values of $Q$ on grid points
have the common denominator $8DA_j^2$ and numerators of bit length polynomial in
$I+j+\mu K_\mu$. The capacities of Lemma~\ref{lit:lem:chaincut} are sums of at most
$O(n^2)$ products of input coefficients with grid coordinates and of correction values, so they
have the same common denominator and numerators of polynomial bit length. After multiplication by
the common denominator, capacities are integers, and the Edmonds--Karp algorithm performs a
number of augmentations polynomial in the graph size, each on integers bounded by the total
capacity. The number of stages per trial is $O(I+q+1)$ by~\eqref{eq:stagebudget}, and the number
of trials is $O(\log\kappa)$. Multiplying these polynomial bounds gives the first claim. The
exact claim follows from the proof of Theorem~\ref{thm:exact}, which uses only the approximation
algorithm, the height bounds of Lemma~\ref{lem:height}, rational reconstruction and exact
checks; it requires $q=\poly(I)+O(\log\kappa)$.
\end{proof}

\begin{corollary}[Balanced after deleting a supplied core]\label{lit:cor:core}
Let $F$, $X$, $L>0$, growth and $\kappa$ be as in Theorem~\ref{lit:thm:balanced}(2), except that
$F$ need not be balanced. Suppose a set $C$ of $k$ coordinates is supplied such that the
quadratic obtained from $F$ by fixing $x_C$ is balanced, that is, there are signs $\sigma$ with
$\sigma_i\sigma_jH_{ij}\le0$ for all $i\neq j$ outside $C$. Then the conclusions of
Theorem~\ref{lit:thm:balanced}(2) hold with work multiplied by
$K^{k}$, where $K=O(\sqrt\kappa\log n)$ bounds the labels per coordinate.
\end{corollary}

\begin{proof}
Fixing the labels of the core coordinates on their grids changes only unary and constant terms of
the remaining quadratic. Hence $\beta$ and each min-marginal are minima, over at most $K^k$ core
label vectors, of balanced grid problems solved by Lemma~\ref{lit:lem:chaincut}; min-marginals of
core coordinates restrict the enumeration. The remainder of the proof of
Theorem~\ref{lit:thm:balanced} is unchanged.
\end{proof}

This removes the restriction to nonpositive residual diagonals in the cut-based recourse of the
report: residual coordinates with positive diagonal are gridded and filtered like the core, and
the residual graph may be dense.

\begin{remark}[Position of Theorem~\ref{lit:thm:balanced}]
After sign changes, a balanced box quadratic is a continuous submodular quadratic. For it,
\cite{BurerNatarajanWillemsen2026v3} prove that a polynomial-size SDP relaxation is exact when
$n\le3$, give a four-variable gap, and state that the complexity for $n\ge4$ is open
(Theorem~1, Example~4 and footnote~2). Exactness of the basic SDP relaxation is known under an
additional sign condition on the linear term \cite{KimKojima2003}, and the case of nonpositive
diagonal is solved by endpoint reduction and a cut, as in part (1)
\cite{BurerNatarajanWillemsen2026v3,Padberg1989}. Discretization methods for continuous
submodular minimization \cite{Bach2019,Bach2018Isotonic,AxelrodLiuSidford2020} are polynomial
in $1/\varepsilon$. Theorem~\ref{lit:thm:balanced} instead has polynomial dependence on
$\log(1/\varepsilon)$ and on $\kappa$. Since $\kappa$ can be exponential in $I$, it does not
settle the open complexity question; it shows that under quadratic growth the cost is
polynomial in the conditioning.
\end{remark}

\section*{C. Optimal sets}

\begin{proposition}[Face criterion for coordinatewise concave quadratics]\label{lit:prop:faces}
Let $F(x)=x^{\mathsf T}Ax+c^{\mathsf T}x+e$ on a mixed box $X$ whose continuous hull is
$[l,u]$, $l<u$, with integer endpoints for integer coordinates, and suppose $A_{ii}\le0$ for all
$i$. Let $\mathcal V=\prod_i\{l_i,u_i\}$, $F^*=\min_{\mathcal V}F$, and for $x\in[l,u]$ let
$\Phi(x)=\prod_i\Phi_i(x_i)$ with $\Phi_i(x_i)=\{x_i\}$ if $x_i\in\{l_i,u_i\}$ and
$\Phi_i(x_i)=\{l_i,u_i\}$ otherwise. Then $F^*=\min_XF$ and, for $x\in X$,
\[
F(x)-F^*=\sum_{v\in\mathcal V}\bigl(F(v)-F^*\bigr)W(v;x)-\sum_iA_{ii}(x_i-l_i)(u_i-x_i),
\qquad W(v;x)=\prod_i\frac{|x_i-\bar v_i|}{u_i-l_i},
\]
where $\bar v_i$ is the endpoint different from $v_i$. Consequently $F(x)=F^*$ if and only if
(a) $x_i\in\{l_i,u_i\}$ whenever $A_{ii}<0$, and (b) $F(v)=F^*$ for every $v\in\Phi(x)$.
If moreover $F=\sum_t\varphi_t(x_{B_t})$ for a rooted tree decomposition, and
$r_t\ge0$ are the endpoint Bellman residuals of the min-sum recursion on $\mathcal V$, so that
$F(v)-F^*=\sum_tr_t(v_{B_t})$ on $\mathcal V$, then (b) holds if and only if
$r_t(w)=0$ for every $t$ and every $w$ in the projection of $\Phi(x)$ to $B_t$.
\end{proposition}

\begin{proof}
Let $Y$ have independent coordinates with $Y_i\in\{l_i,u_i\}$ and $\E Y_i=x_i$; then
$\Pr(Y=v)=W(v;x)$, and $W(v;x)>0$ exactly for $v\in\Phi(x)$. Independence gives
$\E Y_iY_j=x_ix_j$ for $i\neq j$ and $\E Y_i^2=x_i^2+(x_i-l_i)(u_i-x_i)$. Hence
$\E F(Y)=F(x)+\sum_iA_{ii}(x_i-l_i)(u_i-x_i)$, which is the displayed identity after subtracting
$F^*=\sum_vF^*W(v;x)$. Every term on the right is nonnegative, so $F\ge F^*$ on $[l,u]\supseteq
X$, and $\mathcal V\subseteq X$ gives $F^*=\min_XF$. The left side vanishes exactly when every
term vanishes, which is (a) and (b). For the factored statement, (b) says
$\sum_tr_t(v_{B_t})=0$ for all $v\in\Phi(x)$; since each $r_t\ge0$ and $\Phi(x)$ is a product
set, every $w$ in the projection of $\Phi(x)$ to $B_t$ extends to some $v\in\Phi(x)$, which gives
the equivalence. The residual identity is the telescoping of the min-sum recursion: with
$m_t(v_{S_t})=\min_{v_{B_t\setminus S_t}}[\varphi_t+\sum_{c}m_c]$ over endpoint labels and
$m_{\mathrm{root}}=F^*$, put $r_t=\varphi_t+\sum_cm_c-m_t\ge0$; summing over $t$, each nonroot
message appears once with each sign.
\end{proof}

\begin{remark}[Attribution]
For multilinear functions on a box ($A_{ii}=0$), the criterion ``the optimal set is the union of
the faces all of whose vertices are optimal'' is \cite[Lemma and Proposition~1, p.~96]{Rosenberg1972}.
The zero-residual description of all optimal labelings of a tree-structured discrete problem is
the reparameterization (equivalent transformation) view of max-sum problems
\cite[Section~III-D, Theorem~4]{Werner2007}, \cite[Proposition~1]{WainwrightJaakkolaWillsky2005}.
Endpoint reduction for nonpositive diagonals is used in
\cite{DelPiaKhajavirad2026,Khajavirad2026PolyBox}. Proposition~\ref{lit:prop:faces} combines these:
it adds the exclusion (a) for strictly concave coordinates, mixed integer boxes, and a factored
certificate of size $O(N2^p)$.
\end{remark}

\begin{proposition}[Diagonal Lagrangian certificate and the full optimal set]\label{lit:prop:diag}
Let $F(x)=\tfrac12x^{\mathsf T}Hx+b^{\mathsf T}x$ on $B=[l,u]$, $l<u$, and let $s\in B$ satisfy
the box KKT conditions: $\partial_iF(s)\ge0$ if $s_i=l_i$, $\partial_iF(s)\le0$ if $s_i=u_i$, and
$\partial_iF(s)=0$ otherwise. Define $d_i=2\partial_iF(s)/(u_i-l_i)$ if $s_i=l_i$,
$d_i=-2\partial_iF(s)/(u_i-l_i)$ if $s_i=u_i$, $d_i=0$ otherwise, and $D=\operatorname{diag}(d)$.
\begin{enumerate}
\item For every $x\in\R^n$,
$F(x)-F(s)=\tfrac12(x-s)^{\mathsf T}(H+D)(x-s)+\tfrac12\sum_id_i(x_i-l_i)(u_i-x_i)$, and $d\ge0$.
\item If $H+D\succeq0$, then $s$ is a global minimizer on $B$ and
\[
\arg\min_BF=\{x\in B:\ (H+D)(x-s)=0,\ d_i(x_i-l_i)(u_i-x_i)=0\ \text{for all }i\}.
\]
\item Under (2), every $t\in\arg\min_BF$ satisfies the box KKT conditions and defines the same
vector $d$. Thus the test $H+D\succeq0$ is passed at every global minimizer as soon as it is
passed at one.
\end{enumerate}
\end{proposition}

\begin{proof}
(1) Nonnegativity of $d$ is the KKT sign condition. Exact expansion gives
$F(x)-F(s)=\nabla F(s)^{\mathsf T}(x-s)+\tfrac12(x-s)^{\mathsf T}H(x-s)$. For $s_i=l_i$,
$\tfrac12d_i[(x_i-l_i)^2+(x_i-l_i)(u_i-x_i)]=\tfrac12d_i(x_i-l_i)(u_i-l_i)=\partial_iF(s)(x_i-s_i)$;
for $s_i=u_i$ the same computation gives $\tfrac12d_i(u_i-x_i)(u_i-l_i)=\partial_iF(s)(x_i-s_i)$;
for interior $s_i$ both sides are zero. Summing proves the identity.
(2) On $B$ both terms are nonnegative, so $F\ge F(s)$, and $F(x)=F(s)$ exactly when both vanish;
for a positive semidefinite matrix the quadratic form vanishes exactly on its kernel.
(3) Let $t$ be optimal. By (2), $(H+D)(t-s)=0$, so $\nabla F(t)=\nabla F(s)+H(t-s)=\nabla F(s)-D(t-s)$,
and $t_i\in\{l_i,u_i\}$ whenever $d_i>0$. If $d_i>0$ and $t_i=s_i$, then
$\partial_iF(t)=\partial_iF(s)$. If $d_i>0$ and $t_i$ is the other endpoint, say $s_i=l_i$ and
$t_i=u_i$, then $\partial_iF(t)=d_i(u_i-l_i)/2-d_i(u_i-l_i)=-d_i(u_i-l_i)/2\le0$, which is the
KKT sign at an upper bound and yields the same $d_i$; the case $s_i=u_i$, $t_i=l_i$ is symmetric.
If $d_i=0$, then $\partial_iF(s)=0$ by the definition of $d$ and the KKT conditions, so
$\partial_iF(t)=\partial_iF(s)-d_i(t_i-s_i)=0$ and the $i$th entry for $t$ is $0$ whatever the
position of $t_i$. Hence $t$ is a KKT point with the same $d$.
\end{proof}

\begin{remark}[Attribution]
With $\lambda_i=d_i/2$, (1) states that
$F(x)+\sum_i\lambda_i(x_i-l_i)(x_i-u_i)=F(s)+\tfrac12(x-s)^{\mathsf T}(H+D)(x-s)$, the Lagrangian
of the quadratic description $(x_i-l_i)(x_i-u_i)\le0$ of the box. Statement (2) is the box
Lagrangian sufficient condition stated as \cite[Corollary~2]{LiWuQuan2015}; related earlier
conditions are \cite[Theorem~2.3]{BeckTeboulle2000} for $\{-1,1\}^n$ and
\cite{JeyakumarRubinovWu2006} for boxes (see \cite[Proposition~2.1]{HuyJeyakumarLee2006}).
The description of the whole optimal set and statement (3) are instances of the classical fact
that a multiplier certifying a zero duality gap is a multiplier for every optimal solution;
for the convexified Lagrangian this is Lemma~\ref{lit:lem:mangasarian}
\cite{Mangasarian1988,BurkeFerris1991,JeyakumarLeeDinh2004}. The direct proof above is given
because the paper needs exactly this form.
\end{remark}

\section*{D. Mean-preserving rounding in totally unimodular fibres}

\begin{lemma}[Corner rounding in a dilated TU fibre]\label{lit:lem:tu}
Let $A\in\Z^{m\times n}$ be totally unimodular, $r\in\Z^m$, $\alpha,\omega\in\Z^n$, $h>0$, and
$k\in\Z^n$. Let $P=\{t\in[0,1]^n:\ At\le r,\ \alpha\le t\le\omega\}$. Then every $t\in P$ is a
convex combination of points of $P\cap\{0,1\}^n$. Consequently, if $x=h(k+t)$ with $t\in P$,
there is a random vector $Y$ with values in $h(k+(P\cap\{0,1\}^n))$, $\E Y=x$, and
$\E(Y_i-x_i)^2=h^2t_i(1-t_i)\le h^2/4$ for each $i$. If $\nabla^2F\preceq LI$ on the convex hull
of the support of $Y$ and $x$, then $\E F(Y)\le F(x)+Lnh^2/8$. Moreover, if $x_i$ lies in the
grid interval $[hk_i,h(k_i+1)]$, then $Y_i$ takes only its two endpoints.
\end{lemma}

\begin{proof}
The matrix obtained by stacking $A$, $I$ and $-I$ is totally unimodular, and the right-hand
sides $r$, $\min\{\omega,1\}$ and $-\max\{\alpha,0\}$ are integral, so $P$ is an integral
polytope \cite{HoffmanKruskal1956}; its vertices lie in $[0,1]^n\cap\Z^n=\{0,1\}^n$. Write
$t=\sum_v\pi_vv$ as a convex combination of vertices and let $\Pr(Y=h(k+v))=\pi_v$. Then
$\E Y=x$, and since $Y_i\in\{hk_i,h(k_i+1)\}$ with mean $x_i$, its variance is
$h^2t_i(1-t_i)$. Taylor's formula with the Hessian bound gives
$F(Y)\le F(x)+\nabla F(x)^{\mathsf T}(Y-x)+\tfrac L2\norm{Y-x}^2$; taking expectations removes
the linear term and gives the bound.
\end{proof}

\begin{remark}
In the TU extension of the paper, $P$ is the fibre $\{t: At\le2^j(b-Bz)-Ak\}$ of the cell
containing a feasible point at mesh $h=2^{-j}$, with the box bounds rescaled; integrality of
$l,u,B,b$ and the dyadic mesh make all right-hand sides integral. Only the integrality of $P$ is
used; any Carath\'eodory decomposition provides the rounding. The rounding is correlated, which is
why a full Hessian bound rather than a diagonal one is needed. Marginal-preserving rounding under
hard constraints is classical in approximation algorithms
\cite{GandhiEtAl2006,ChekuriVondrakZenklusen2010}; the requirement here is only exact
feasibility and the mean.
\end{remark}

\section*{E. Draft: related work for the extensions}

\paragraph{Value functions, multiparametric programming and condensation.}
Eliminating private variables through their optimal response is a standard form of
decomposition: Benders and generalized Benders decomposition optimize over value functions of
subproblems \cite{Benders1962,Geoffrion1972}, partial minimization of a quadratic is a Schur
complement \cite{Zhang2005Schur}, and condensing eliminates states in model predictive control
\cite{BockPlitt1984,JerezKerriganConstantinides2012,Axehill2015}. When the eliminated block has
active bound constraints, multiparametric quadratic programming gives piecewise-affine optimizers
and piecewise-quadratic values on polyhedral critical regions
\cite{BemporadEtAl2002,TondelJohansenBemporad2003}, also for positive semidefinite Hessians with
nonunique optimizers \cite{SpjotvoldTondelJohansen2005,SpjotvoldTondelJohansen2007,
SpjotvoldEtAl2006,PatrinosSarimveis2011}. Value functions of parameters that enter the objective
linearly over a fixed feasible set are concave \cite{MangasarianRosen1964,FiaccoKyparisis1986},
and infima preserve semiconcavity \cite{CannarsaSinestrari2004}. Our use of these facts is
different in purpose: we need an \emph{upper coordinate curvature} bound for the reduced
objective, valid globally, so that the corrected grid remains a certified lower bound. The
closest algorithmic precedents eliminate continuous components exactly: Khajavirad
\cite{Khajavirad2026PolyBox} and Del Pia and Khajavirad \cite{DelPiaKhajavirad2026} reduce box
QPs to binary problems on small neighborhoods of positive-diagonal components, evaluating each
component's value function at all binary neighbor labels, and use the concave kinks of univariate
value functions for forests. In our setting the retained neighbors are continuous, gridded and
filtered, and the parameter is curvature over growth. Low negative inertia with convex recourse
is treated by spatial branching in \cite{Vavasis1992,LuoBaiLimPeng2019}.

\paragraph{Cut oracles, submodularity and smoothing.}
Binary quadratic functions with nonpositive pair coefficients are minimized by one minimum cut
\cite{Ivanescu1965,PicardRatliff1975,BorosHammer2002,KolmogorovZabih2004}, after sign switching
whenever the signed interaction graph is balanced \cite{Harary1953,HammerHansenSimeone1984,
Padberg1989}. Variables with nonpositive diagonal can be fixed to endpoints
\cite{Rosenberg1972,DelPiaKhajavirad2026,Khajavirad2026PolyBox}, so continuous quadratics that
are submodular and concave in each coordinate reduce to such cuts
\cite{BurerNatarajanWillemsen2026v3}. Pairwise problems that are submodular with respect to label
orders reduce to cuts on threshold encodings \cite{SchlesingerFlach2006,Ishikawa2003,Hochbaum2001}.
For general continuous submodular quadratics the semidefinite relaxation of
\cite{BurerNatarajanWillemsen2026v3} is exact only up to dimension three, while
\cite{KimKojima2003} cover a sign-restricted case and discretization methods
\cite{Bach2019,Bach2018Isotonic,AxelrodLiuSidford2020} are polynomial in the inverse accuracy.
Partial minimization preserves submodularity \cite{Topkis1978,Topkis1998}; this underlies the
mixed continuous--discrete methods of \cite{GomezHan2025,BuntonTabuada2022} and our mixed
recourse oracle, whose certificates are convex combinations of greedy extreme bases
\cite{Edmonds1970,Cunningham1985,IwataFleischerFujishige2001,Schrijver2000}. The perturbation
model for the core is in the tradition of smoothed analysis
\cite{SpielmanTeng2004,BeierVocking2006,RoglinVocking2007}; discrete uniform weights also appear in
the isolation lemma \cite{MulmuleyVaziraniVazirani1987}, and low-dimensional nonconvexity with
random perturbations is studied in \cite{KelnerNikolova2007}.

\paragraph{Constraints.}
Integrality of dilated totally unimodular fibres is the Hoffman--Kruskal theorem
\cite{HoffmanKruskal1956,Schrijver1986}; box-integrality is analyzed in
\cite{ChervetGrappeRobert2021}. Proximity between continuous and integer optima of convex
problems \cite{CookGerardsSchrijverTardos1986,GranotSkorinKapov1990,HochbaumShanthikumar1990}
addresses a different question: we round arbitrary feasible points of a nonconvex objective and
pay a curvature term. Bounded-width constrained polynomial optimization admits polynomial-size
linear programming approximations with scaled feasibility tolerance \cite{BienstockMunoz2018};
our TU extension is narrower but exactly feasible.

\paragraph{Optimal sets and certificates.}
The faces of optimal points of multilinear box problems were characterized by Rosenberg
\cite{Rosenberg1972}; tree reparameterizations describe all optimal labelings of discrete
problems \cite{WainwrightJaakkolaWillsky2005,Werner2007}. Diagonal Lagrangian certificates for
binary and box quadratics are given in \cite{BeckTeboulle2000,JeyakumarRubinovWu2006,Pinar2004,
LiWuQuan2015}, and solution sets of convex programs are characterized by constant gradients
\cite{Mangasarian1988,BurkeFerris1991,JeyakumarLeeDinh2004}. Quadratic growth toward the optimal
set of a quadratic program on a polytope follows from the error bound of Luo and Sturm
\cite{LuoSturm2000}, without a computable constant.

\paragraph{Active faces.}
Sign tests on gradient enclosures are the monotonicity test of interval global optimization
\cite{Hansen1980,HansenWalster2004} and of interval constraint propagation
\cite{ArayaTrombettoniNeveu2010}; Bernstein expansions give polynomial range bounds
\cite{Garloff1986}. Identification of active constraints under strict complementarity or error
bounds is classical in local optimization \cite{BurkeMore1988,Wright1993,FacchineiFischerKanzow1998}.
Our contribution there is only the explicit precision term for reaching such a region from a
certified global enclosure.

\paragraph{Structural limits.}
Del Pia and Khajavirad prove strong NP-hardness at treewidth two with bounded coefficients
\cite{DelPiaKhajavirad2026}; sparse second-order cone and semidefinite formulations, the latter with RLT products, are exact
under graph conditions \cite{DeyKhajavirad2026,Khajavirad2026SparseSDP}. Chordal sparsity guarantees
positive semidefinite completions \cite{GroneEtAl1984,VandenbergheAndersen2015}, and sparse
moment relaxations exploit chordal structure and converge under a running-intersection property
\cite{Lasserre2001,WakiEtAl2006,Lasserre2006,Laurent2009}; locally consistent moments need not
come from a joint measure, the continuous counterpart of the gap between local and global
marginal polytopes \cite{SheraliAdams1990,WainwrightJordan2008}. Complete parametric descriptions
can have exponentially many pieces \cite{Murty1980,MairalYu2012,GartnerJaggiMaria2012}.
